\section{Related Work}
\label{sec:related}

\paragraph{Prompt compression.}
LLMLingua~\citep{jiang2023llmlingua} and LLMLingua-2~\citep{pan2024llmlingua2}
show that prompts can be compressed 3--20$\times$ with minimal quality loss,
demonstrating substantial token redundancy.
\citet{chevalier2023autocompressors} train models to summarise prompts into
soft tokens.
More recent pipelines such as CompactPrompt~\citep{wang2025compactprompt}
combine token pruning with n-gram abbreviation for end-to-end compression.
These methods operate \emph{at inference time, per instance}---compressing a
written prompt by removing low-information tokens.
Our work operates \emph{at authoring time, per task}: we provide priors on
which prompt mechanisms yield quality gains for a given task type, before
any specific instance is constructed.
The two approaches are complementary: author the right mechanisms (our work),
then compress the wording within them (LLMLingua).

\paragraph{Prompt engineering.}
Chain-of-thought~\citep{wei2022cot} and few-shot~\citep{brown2020gpt3}
prompting add tokens in exchange for quality gains on reasoning tasks.
\citet{zhou2022large} show that task-description quality matters more than
length for many tasks.
Recent surveys of prompting strategies~\citep{sahoo2024systematic,
schulhoff2024prompt} and automatic prompt
optimisation~\citep{amatriain2025survey} catalogue dozens of techniques but
do not measure their marginal effect as a function of token count.
Our work complements these by quantifying the return of each successive
elaboration layer and showing that marginal returns vary systematically
across task types.

\paragraph{Long-context performance.}
\citet{liu2023lost} demonstrate that LLMs underutilise information positioned
in the middle of long contexts, a degradation effect distinct from our
finding, which concerns the marginal benefit of additional \emph{instruction}
tokens before any degradation.
\citet{levy2024same} and \citet{hsieh2024ruler} study how performance varies
with context length on retrieval and reasoning benchmarks, finding similar
non-monotone effects at much larger scales.

\paragraph{Prompt sensitivity and robustness.}
\citet{webson2022prompt} and \citet{mizrahi2024state} show that seemingly
small prompt changes cause large quality swings.
Our design controls for this by holding content constant and varying only
elaboration depth, isolating the information-content signal from phrasing
effects.
The randomised ablation further controls for ordering effects.

\paragraph{Prompt sensitivity: formatting axis.}
Several recent works study how LLM performance varies under prompt
perturbations.
\citet{zhuo2024prosa} introduce ProSA, a metric quantifying sensitivity
to paraphrased prompts while holding information constant.
\citet{sclar2024quantifying} show that minor formatting choices
(separators, casing) substantially affect in-context learning accuracy.
\citet{voronov2024mind} demonstrate that inconsistent prompt formatting
confounds evaluation of in-context learning improvements.
These works study sensitivity along the \emph{formatting} axis---varying
phrasing or layout while holding information constant.
Our work studies the complementary \emph{information} axis: we hold
phrasing fixed (via strict additive nesting) and vary how much
task-relevant information the prompt contains.

\paragraph{Scaling laws.}
Neural scaling laws~\citep{kaplan2020scaling,hoffmann2022chinchilla} describe
how model quality grows with compute and training data.
We study an analogous phenomenon at the \emph{prompt level}: how quality
grows with instruction tokens, and where this growth saturates.
To our knowledge, no prior work has fitted and statistically tested saturation
curves across multiple tasks, models, and prompt orderings in a controlled
additive design.

\paragraph{LLM-as-judge evaluation.}
\citet{zheng2023judging} show that strong LLM judges match human preferences
at over 80\% agreement, establishing the viability of automated evaluation.
Recent work has shown that different LLM judges exhibit systematically
different score calibrations~\citep{chen2025scoring,li2026same}, motivating
our dual-judge design with both gpt-4o-mini and gemini-2.0-flash.
